\chapter{Vertical slices}
\label{ch:slices}

\section*{What this chapter is for}

How to choose what to build first, and why the instinct most engineers have
\dash{} build the foundations properly, then build on them \dash{} produces the
worst possible outcome under a clock.

\section{Horizontal loses}

The horizontal approach builds a layer at a time: the schema, then the data
access, then the services, then the interface. It is how you would build
something real, it produces better foundations, and in a timed exercise it
produces nothing that runs.

The arithmetic is brutal. Four layers, and you are seventy per cent through
each when the time goes. Seventy per cent of a schema is nothing. Seventy per
cent of four layers is nothing. There is no demonstration, no account of what
works, and nothing to talk about except intent.

A vertical slice is the thinnest complete path through every layer at once.
One entity, one operation, one screen, one assertion. It is ugly, it is
incomplete, and it \emph{runs}.

\begin{kitnote}
The property that matters is not that it is smaller. It is that it is
\emph{complete}. A complete narrow thing tells you that your assumptions about
every layer were right, which is information you cannot get any other way and
which you want as early as possible. A wide incomplete thing tells you nothing
about whether it would have worked.
\end{kitnote}

\section{Choosing the first slice}

\begin{kitmethod}
\textbf{It must be the thing discovery said has to work.} Question two of the
six. If the answer was ``a customer can place an order'', the first slice is a
customer placing an order, not a catalogue.

\textbf{It must go through every layer.} Store, logic, boundary, whatever the
user touches. If one layer is missing, you have not proved the path.

\textbf{It must contain the one hard thing.} The conflict case, the external
dependency, the operation that is irreversible. Chapter~\ref{ch:design} phase
7. Putting the hard thing in slice one is the whole technique: it is the part
that can fail in ways you did not predict, and you want it failing at minute
forty rather than minute ninety.

\textbf{It must be demonstrable.} You can show somebody, in the product, that
it works.
\end{kitmethod}

\begin{kittrap}
The tempting first slice is the one you are confident about \dash{} the
catalogue, the list view, the read path. It goes quickly, it produces something
on screen, and it feels like progress.

It defers every risk to the second half of the session, where there is no room
left. The hard thing does not get easier by being postponed; it gets more
expensive, because by then it has to fit into decisions you have already made
around it.
\end{kittrap}

\section{What a slice is allowed to skip}

It has to be complete through the layers. It does not have to be complete in
any other sense, and being explicit about that at the time is what keeps it
from reading as carelessness.

\begin{kittable}{@{}lX@{}}
\textbf{Allowed to skip} & \textbf{Because} \\
\midrule
Authentication & a stub identity, named as a stub, costs one line and proves nothing was skipped by accident \\
Validation beyond the path & the path is what is being proved \\
Styling & nobody is assessing it \\
Configuration & a constant with a comment is a decision; a configuration system is a project \\
Every entity but one & the second entity is slice two \\
Error handling beyond the one case & the one case is the one that matters \\
\end{kittable}

\textbf{Not allowed to skip:} the assertion. A slice that runs but is not
checked has not proved anything, and in a round where the criteria include how
you validate your own work, it is the omission that gets noticed.

\judgement{One test on the hard thing is worth more than ten on the easy parts,
and writing it in front of somebody is one of the highest-value two minutes
available in the whole session.}

\section{The seam}

Every slice should leave an obvious place for the next thing to attach, and
saying where that is, out loud, is what turns a narrow slice from a limitation
into a plan.

``The second entity attaches here, and the only change is another row in this
mapping'' is a sentence that answers the unasked question \dash{} whether you
have built something that can grow or something that happens to work.

\section{Slice two, and the decision not to start it}

If slice one finishes early, slice two is either the second-most-important
thing or the hardening of the first. Both are defensible; choose deliberately
and say which.

If slice one finishes late, do not start slice two. The run sheet already made
this decision. \judgement{Two half-finished slices demonstrate less than one
finished one and cost the demonstration, which is the only artefact the round
actually produces.}

\begin{drill}
Take a system you have built and write down what its first vertical slice
would have been: the one path, through every layer, containing the hardest
thing. Then check the last part honestly. Most people's first slice avoids the
hard thing, and finding out that your instinct does that is the point of the
exercise.
\end{drill}

\section*{What this chapter does not claim}

Vertical slicing is not novel and is not this book's idea. What is specific
here is the instruction to put the hardest thing in the first slice, which
trades a smoother start for earlier information, and is a judgement about where
risk should sit when the clock is the binding constraint.

Under different constraints \dash{} a week rather than two hours \dash{} the
horizontal approach is often correct. The claim is about timed exercises.
